\section{Threat Surface and System Context}
\label{sec:system-context}

The method is placed wherever externally influenced media becomes eligible for
local use. The upstream source may be an authenticated message, mail attachment,
cloud-synchronized object, browser download, removable medium, or direct import.
The downstream consumer may be durable storage, a preview generator, a media
decoder, an export interface, or another application. Upstream provenance does
not substitute for validation of the exact octet string crossing this local
trust boundary.

\begin{figure*}[t]
\centering
\resizebox{0.98\textwidth}{!}{%
\begin{tikzpicture}[
  box/.style={draw,rounded corners=2pt,align=center,font=\small,inner sep=3pt},
  source/.style={box,fill=blue!7,minimum width=42mm,minimum height=9mm},
  boundary/.style={box,fill=gray!12,minimum width=72mm,minimum height=10mm},
  threat/.style={box,fill=red!7,minimum width=43mm,minimum height=9mm},
  control/.style={box,fill=green!8,minimum width=43mm,minimum height=9mm},
  surface/.style={draw,rounded corners=2pt,dashed,inner sep=4pt},
  arrow/.style={-{Latex[length=2mm]},thick},
  enter/.style={-{Latex[length=2mm]},thick,dashed}
]
\node[source] (s1) at (0,2.5) {messaging after local\\decryption};
\node[source] (s2) at (5.0,2.5) {email, web, and\\cloud synchronization};
\node[source] (s3) at (10.0,2.5) {local or removable\\media import};
\node[boundary] (input) at (5.0,1.2) {untrusted evidence: bytes, portable name,
declared type\\trusted context: policy, deadline, cancellation};
\draw[arrow] (s1) -- (input);
\draw[arrow] (s2) -- (input);
\draw[arrow] (s3) -- (input);

\node[threat] (t1) at (0,-0.2) {evidence conflict\\or type confusion};
\node[threat] (t2) at (5.0,-0.2) {malformed boundary,\\suffix, or polyglot};
\node[threat] (t3) at (10.0,-0.2) {forbidden metadata\\or nested object};
\node[control] (c1) at (0,-1.35) {one admissible profile\\or reject};
\node[control] (c2) at (5.0,-1.35) {authorizing parse and\\output evidence convergence};
\node[control] (c3) at (10.0,-1.35) {allowlisted structural\\reconstruction};

\node[threat] (t4) at (0,-2.9) {retained codec payload\\or decoder trigger};
\node[threat] (t5) at (5.0,-2.9) {size, count, duration,\\or expansion abuse};
\node[threat] (t6) at (10.0,-2.9) {fault, race, overwrite,\\or partial output};
\node[control] (c4) at (0,-4.05) {semantic rebuild\\or block};
\node[control] (c5) at (5.0,-4.05) {checked bounds, deadline,\\and cancellation};
\node[control] (c6) at (10.0,-4.05) {empty failure result and\\non-overwriting commit};

\foreach \i in {1,2,3,4,5,6} {\draw[arrow] (t\i) -- (c\i);}
\node[surface,fit=(t1)(t2)(t3)(c1)(c2)(c3)(t4)(t5)(t6)(c4)(c5)(c6)]
  (surface) {};
\draw[enter] (input.south) -- (surface.north);
\end{tikzpicture}%
}
\caption{Source-independent threat-control map. Upstream channel properties do
not authorize media bytes. Each measurable carrier or publication threat is
assigned a distinct fail-closed control.}
\label{fig:component-architecture}
\end{figure*}

In an end-to-end encrypted application, this boundary lies after authenticated
decryption and before local storage or rendering. The relay can remain blind to
plaintext and the local verdict. In other deployments the same method begins
after download or import. It neither establishes transport confidentiality nor
changes the authentication model of the source channel.

The protected assets are parser and client availability, privacy-sensitive
metadata, integrity of reconstructed content, policy consistency, and integrity
of the final publication transition. Endpoint execution integrity is an asset
of the enclosing deployment but is not guaranteed against a compromised
external decoder by this component alone. The adversary controls the full input
byte string, display name, supplied media
type, internal sizes and offsets, metadata, compressed payload, dimensions,
counts, duration fields, and submission timing. The adversary may construct
malformed files, prefix and suffix camouflage, whole-file polyglots, nested
objects, decompression amplification, decoder regressions, and cancellation
races. Repeated submissions and adaptive selection based on visible success or
failure are permitted. The policy snapshot is trusted. Caller cancellation and
earlier deadlines may shorten work but cannot enlarge the authorized policy or
the host processing budget.

The in-scope threat classes are covered by
\begin{equation}
  \mathcal{T}_{\mathrm{in}}=
  \mathcal{T}_{\mathrm{id}}\cup
  \mathcal{T}_{\mathrm{grammar}}\cup
  \mathcal{T}_{\mathrm{payload}}\cup
  \mathcal{T}_{\mathrm{resource}}\cup
  \mathcal{T}_{\mathrm{publication}},
  \label{eq:threat-cover}
\end{equation}
where the classes denote evidence and type confusion, malformed or ambiguous
whole-file structure, retained or forbidden payload structure, computational
amplification, and fault-driven publication effects. The classes may overlap.
A test record may carry secondary tags, while its prespecified primary class
determines the aggregate denominator and primary falsification oracle.
Arbitrary decoder compromise is a conditional deployment threat because process
supervision alone does not isolate same-user host effects. Harmful decoded
meaning and steganography are content-analysis problems outside
\(\mathcal{T}_{\mathrm{in}}\).

\begin{table*}[t]
\centering
\caption{Adversary capabilities and the corresponding enforced boundary.}
\label{tab:adversary-capabilities}
\small
\begin{tabularx}{\textwidth}{L{0.23\textwidth}YY}
\toprule
Control & Security-relevant effect & Required response \\
\midrule
Name and supplied type & Misrouting, hidden suffix, or type confusion & Portable-name grammar and evidence convergence before dispatch \\
Arbitrary byte layout & Prefix success, trailing object, overlap, or malformed boundary & Offset-zero complete parsing and exact terminal offset \\
Metadata and container structure & Privacy leakage, nested active object, or stale reference & Location-aware allowlist and reconstruction \\
Compressed payload & Exploitation during source decoding or retention of a source representation in the result & Isolate or trust the decoder for transformation-time risk. Use semantic reconstruction when source-payload retention is prohibited \\
Counts, dimensions, and timing & Allocation, expansion, or execution exhaustion & Checked arithmetic, limits, deadline, and cancellation \\
Fault timing and destination state & Partial publication or overwrite & Temporary output, validation before commit, and non-overwriting persistence with deployment-verified atomic visibility \\
\bottomrule
\end{tabularx}
\end{table*}

Malware labels organize scenario provenance but do not enter the output verdict.
The behavioral distinctions defined in \cref{sec:background} prevent a carrier
property from being reported as evidence of replication, delivery role, or
payload objective. The empirical corpus can therefore contain inert carriers
shaped around those mechanisms without implying a malware-family detection rate.

The trusted computing base includes the authenticated application binary, the
loaded immutable policy snapshot, the registered authorizing parsers, trusted
provenance records, and the operating-system enforcement of process and
filesystem primitives. Collision-resistant digests identify and integrity-check
recorded artifacts but do not authenticate their provenance by themselves. The
endpoint is assumed uncompromised before processing begins. Memory-safety
guarantees reduce one class of implementation defects but do not prove parser
correctness. Foreign code and external media tools remain potentially faulty.
Their output is not deliverable until an authorizing parser accepts the complete
object. Process supervision bounds selected liveness, output, and publication
effects. It does not confine arbitrary same-user filesystem or host effects
after decoder code execution. Deployments that include decoder compromise in
scope must add a platform sandbox such as a dedicated account, namespace and
system-call policy, container or mandatory-access-control profile, or
application sandbox.

The method assigns five responsibilities to distinct stages: evidence
convergence and routing, profile-specific reconstruction, complete validation
and rebuilt-output scanning, process supervision, and final result construction
with publication. Candidate-producing stages do not receive authority to
publish their own output. The concrete module and interface realization is
described in \cref{sec:implementation}.

The method assumes an uncompromised endpoint and does not verify the external codec,
operating system, filesystem, or hardware. Its media-only policy excludes active
documents, archives, executables, scripts, PDF, and vector markup. Accepted
visual or auditory content may still be harmful, misleading, age-inappropriate,
or steganographic. Per-request limits also do not replace system-wide admission
control and scheduling.

The eight assurance obligations are complete output consumption (F1), component
closure (F2), output evidence convergence (F3), semantic source-independence
(F4), second-pass stability (F5), bounded execution (F6), validation separate
from construction (F7), and clean-result publication with fault containment
(F8). Each has a separate falsification condition in the evaluation protocol,
and failure in one obligation cannot be offset by success in another.
